\subsection*{S20. Required versus optional inference}

S4 fixed how a partial path is encoded. This scenario asks when a path counts
as partial in the first place.

\begin{lstlisting}[style=json]
{ "a": 1, "b": 2 }     x999
{ "a": 1 }             x1
\end{lstlisting}

\options
\begin{enumerate}[label=\textbf{(\alph*)}]
  \item \textbf{Any absence makes the path partial.} One document out of a
        thousand missing \texttt{b} gives \texttt{b : int option}.
  \item \textbf{A threshold} --- partial only when absent in more than some
        fraction of documents.
  \item \textbf{Any absence makes the path partial, and the observation counts
        are reported} so that a single anomalous document can be found.
        \selected
  \item \textbf{The user declares it} in the configuration, with the corpus
        checked against the declaration.
\end{enumerate}

\decision{A path is partial exactly when it is absent or null in at least one
observation. Inference additionally records, and the interface reports, how
many observations of each kind were seen at every path. The counts are
diagnostic only and never affect what is generated.}

\begin{lstlisting}[style=json]
.a    present with a value  1000 / 1000    -> a : int
.b    present with a value   999 / 1000
      absent                   1 / 1000    -> b : int option
\end{lstlisting}

\paragraph{The report is the natural place for what S4 discards.}
S4 collapses \texttt{null} with absence, so the generated type cannot say
which occurred. The diagnostic can, and should: reporting
present-with-value, explicitly-null and absent as three separate counts
surfaces exactly the information the collapse removes from the type, at no
semantic cost.

\paragraph{Denominators differ by depth.}
The counts are per path, and a path's denominator is the number of
\emph{observations} at that path, not the number of documents. For a root
member that is one per document; for \texttt{.items[].qty} it is one per array
element across the whole input. A ratio of $999/1000$ therefore means
different things at different depths, and the report should say what it is
counting.

\rationale{Option (b) is a threshold on data values, which is the shape of
rule excluded in S9: a field would flip between \texttt{int} and \texttt{int
option} as the corpus grew, making the generated type depend on how much data
happened to be supplied. Option (c) keeps generation purely structural --- it
is (a) semantically, and the generated code is identical --- while supplying
the information a user needs to notice that a field was missing exactly once
and go and look at that document. Given that the tool is used through an
interface where the schema is inspected before the modules are taken away,
that information is worth surfacing next to each field.}

\limitation{A single malformed document makes an otherwise total field
optional for every consumer of the generated module, and the report says so
without fixing it. The only remedy is to remove the document from the input.

The counts are metadata rather than schema, so if they are to be useful
outside the interface they must be exported alongside it, in the same way the
configuration of S9 and S19 must be.}
